% Chapter 0: Abstract

\chapter*{Abstract}

\yvy\ (Indigenous Territory Mapping), \yvytu\ (Chaco Deforestation), \yvyra\ (Carbon Credits), \yrupe\ (Soybean Yield), \kai\ (Wildlife Poaching), and \tatakua\ (Air Quality) — six research threads united by one Python package.

This thesis presents \textbf{SatelliteCV-Paraguay}, an open-source Python toolkit for multi-temporal earth observation of Paraguay. The package unifies six peer-reviewed papers across remote sensing, agriculture, climate, conservation, and social science.

The toolkit leverages pre-trained foundation models (Prithvi, AlphaEarth, DINOv2, LLaVA-1.6), open satellite data (Sentinel-2, Landsat, Planet), and Paraguayan national datasets (paraguay-geodata, MapBiomas Paraguay, INFONA, INDI, SENEPA). All 6 papers share a common pipeline and infrastructure, reducing duplication and enabling efficient execution.

This thesis addresses critical gaps in Paraguay's earth observation capabilities, including: (1) the lack of comprehensive satellite-based land cover change monitoring; (2) limited support for the country's emerging carbon credit market; (3) absence of AI tools that respect indigenous data sovereignty under CARE Principles; (4) limited use of deep learning for wildlife conservation; (5) poor urban air quality forecasting; and (6) lack of precision agriculture tools for soybean yield prediction.

The thesis contributes a reproducible, open-source codebase that enables future Paraguayan researchers to extend the analysis. It provides a foundation for continued earth observation research in Paraguay and serves as a template for similar efforts in other Latin American countries.
